\chapter{Reproducibilidad}\label{anx:reproducibilidad}
\sloppy

\section{Entorno}

Los experimentos se ejecutaron en una estación con GPU NVIDIA RTX 4060 de 8~GB, con Python
3.12 y PyTorch Geometric. Las dependencias exactas están fijadas en \texttt{uv.lock}, en la raíz
del repositorio, y se instalan con \texttt{uv sync}. El conjunto de datos Elliptic se descarga
automáticamente en \texttt{data/} la primera vez.

\section{Corrida principal (sin el cierre del mercado)}

\begin{itemize}
  \item Configuración: \texttt{configs/experiment\_v4\_precierre.yaml}.
  \item Entrenamiento y explicación: \texttt{scripts/v4/run\_precierre.sh} y
        \texttt{scripts/v4/cola\_precierre.sh}. Producen 216 modelos en
        \texttt{results\_models\_v4\_precierre/} y las explicaciones en
        \texttt{results\_v4\_precierre/} y \texttt{results\_v4\_precierre\_grad/}.
  \item Curvas y métricas de rendimiento:
        \texttt{scripts/v4/reunion\_0410/curvas.py --variante precierre}.
  \item Medidas robustas y control de pesos aleatorios:
        \texttt{scripts/v4/cola\_robusta\_precierre.sh}, que ejecuta en secuencia
        \texttt{metrica\_robusta.py} y \texttt{control\_aleatorio.py} con
        \texttt{--variante precierre}.
\end{itemize}

\section{Tablas, figuras y bibliografía del manuscrito}

Las tablas y figuras se regeneran sin GPU desde los resultados guardados:

\begin{itemize}
  \item \texttt{tesis\_v4/scripts/analisis.py}: tablas \LaTeX{} y pruebas estadísticas
        (\texttt{tesis\_v4/tablas/}).
  \item \texttt{tesis\_v4/scripts/figuras.py}: figuras en PDF vectorial
        (\texttt{tesis\_v4/figuras/}).
  \item \texttt{tesis\_v4/scripts/verificar\_referencias.py}: genera
        \texttt{bibliografia.bib} consultando Crossref por DOI y rechaza todo lo que no sea un
        artículo de revista o tenga una retractación registrada. Las únicas excepciones son
        los artículos originales de los métodos y del conjunto de datos usados, listados en
        \texttt{tesis\_v4/referencias/excepciones.bib} con la URL donde se verificaron.
\end{itemize}

\section{Determinismo}

Las semillas de datos (2026), de modelo (42, 43 y 44), de selección de transacciones (1234) y de
los explicadores están fijadas. Aun así, el entrenamiento en GPU no es determinista a nivel de
pesos, porque las agregaciones de mensajes usan sumas atómicas. Repetir la corrida reproduce
las conclusiones, pero no necesariamente todos los decimales.
